\icode{E-MI-SI-L2-261104-002}
如图所示，\(M\)、\(N\)为装在水平面上的两块间距可以调节的光滑竖直挡板，两板间叠放着\(A\)、\(B\)两个光滑圆柱体，现将	两板间距调小些，这时与原来相比，下述结论中正确的是\paren[]

\begin{tikzpicture}[scale=.8,font=\scriptsize]
	\coordinate (A) at($(1,1)+(53:2)$);
	\coordinate (C) at($(A|-0,0)+(1,0)$);
	\draw[fill=orange!60] (0,0) rectangle(-.3,3.6)node[left,black]{\(M\)} (C) rectangle($(C)+(.3,3.6)$)node[right,black]{\(N\)};
	\draw (1,1) circle(1) (A) circle(1);
	\fill (1,1)node[right]{\(B\)} circle(.05) (A)node[right]{\(A\)} circle(.05);
	\foreach \x in {-.8,-.6,...,3.8}
		\draw[gray] (\x,-.2)--++(.2,.2);
	\draw[thick] (-.8,0)--(4,0);
\end{tikzpicture}
\begin{choices}
	\item\(N\)板对圆柱体\(A\)的弹力变小
	\item 圆柱体\(A\)对圆柱体\(B\)的弹力变大
	\item 水平面对圆柱体\(B\)的弹力变大
	\item 水平面对圆柱体\(B\)的弹力变小
\end{choices}
\begin{answer}
	A\qquad 提示：受力分析，并构建力的矢量三角形
\end{answer}